% New English translation from checked cote 25, site 42 (PDF page 43).
% AI-translated by GPT-6 Astra (OpenAI, ChatGPT Pro mode).
\scanpage{25/42}
% NOTE: The raised supsi belongs to the cancelled Y as a superscript; the cancellation crosses the base expression. All diagram arrows are penned. A separate dense blot beside the upper0 has no secure character value.
\mspage{\hand{38}}

$\underline{Y}^{\mathrm{sing}} = P(\underline{\hand{\nu}}_{X/P}) \hand{\longrightarrow} \overset{\hand{\vee}}{P}$. Then H' is “orthogonal” to x, and in particular
if dim H' = r-1 (which by b) is equivalent to the fact that \var{\xout{yx/x$\underline{Y}^{\unsure{supsi}}$xxxxx}}{}{typewriter strike-out}{ega25_25_042-v0001}
y is a \var{\inserted[typed]{geom.}}{geom.}{insertion}{ega25_25_042-v0002} \uline{ordinary} singular point on $Y_{\hand{\xi}}$), x is uniquely determined
by H'.

Proof. We have already proved a), and b) is a consequence of it
by the developments in par.16. It remains to prove c), which is
again a consequence of an essentially trivial compatibility
relation. I merely copy here, in the special case under
consideration (but with arbitrary S), diagram (D') of par.16, which
takes the following form (where we write $\underline{Y}^1$ for $\underline{Y}^{\mathrm{sing}}$):

\diagram{
\begin{tikzcd}[column sep=small,ampersand replacement=\&,nodes in empty cells]
 \& \& 0 \& \& \\
 \& \& \hand{\underline{\Omega}}^1_{X/S\,\underline{Y}^1} \arrow[u] \& \& \\
0 \arrow[r] \& \underline{P}^1(\underline{O}(1,1))\otimes\underline{O}_{\underline{Y}^1} \arrow[r] \arrow[ur, "\hand{\unsure{\nu}.\mu}"] \& (\hand{\underline{\Omega}}^1_{X/S})_{\underline{Y}^1} + (\hand{\underline{\Omega}}^1_{\overset{\hand{\vee}}{P}/S})_{\underline{Y}^1} \arrow[r] \arrow[u] \& \hand{\underline{\Omega}}^1_{\underline{Y}^1/S} \arrow[r] \& 0 \\
 \& \underline{O}_{\underline{Y}^1}(-1,-1) \arrow[u] \arrow[r, "\hand{u}"] \& \hand{\underline{\Omega}}^1_{\overset{\hand{\vee}}{P}/S\,\underline{Y}^1} \arrow[u] \arrow[ur, "\hand{\nu}"'] \& \& \\
\hand{\underline{\omega}}^{-2}_{X/Y}\otimes\underline{O}_{\underline{Y}^1}(-2,-2) \arrow[ur] \& \& \hand{0} \arrow[u] \& \&
\end{tikzcd}
}

The orthogonality relation c) simply expresses that the composite
$\hand{\nu}u$ in this diagram is zero; I dispense with developing
this second exercise in syntax any further. I suppose people were
used to choosing coordinates, and perhaps they saw things more
clearly that way ....

\uline{Corollary} 8.\var{\over{\illegible{1}}{8}}{8}{writing over}{ega25_25_042-v0003}. Suppose S=Spec(k), \var{\xout{\illegible{1} so that $\underline{Y}^{\mathrm{sing}}$}}{}{typewriter strike-out}{ega25_25_042-v0004} \var{\inserted{$f : X \to P$ an immersion}}{$f : X \to P$ an immersion}{insertion}{ega25_25_042-v0005} and X
irreducible \var{\inserted{(still smooth over k)}}{(still smooth over k)}{insertion}{ega25_25_042-v0006}. \var{\xout{\unsure{Then}}}{}{typewriter strike-out}{ega25_25_042-v0007} Then $\underline{Y}^{\mathrm{sing}}$ is smooth and irreducible; let T be
\tnote{25-042-1}{The diagram is carried over literally, including its lower-left omega expression and the plus sign in its middle object. The source spellings quelconqque and symtaxe are translated arbitrary and syntax. Tjrs is translated still in the inserted hypothesis.}

\var{\margin{\hand{42}}}{\placedmargin{\hand{42}}}{marginal note placed}{ega25_25_042-v0008}
